\section*{Bab \ref{ch:g12:phylogenetic-trees} --- Membaca Kekerabatan: Pohon Filogenetik}

\begin{solution}{exo:g12:phylogenetic-trees:1}
Ujung: \omterm{def:g10:biodiversity-scales:species}{spesies} yang dibandingkan. Simpul: nenek moyang bersama terkini
mereka yang bersifat hipotetis. Akar: nenek moyang bersama semuanya.
\end{solution}

\begin{solution}{exo:g12:phylogenetic-trees:2}
Kerabat terdekat kadal adalah mencit dan manusia sekaligus (keduanya
berbagi simpul tepat di bawah cabang kadal). Kerabat terdekat katak
adalah seluruh kelompok bertiga lainnya, sama dekatnya.
\end{solution}

\begin{solution}{exo:g12:phylogenetic-trees:3}
Keadaan yang muncul pada seorang nenek moyang dan diwarisi oleh
keturunannya; berbagi keadaan itu berarti berbagi nenek moyang
tersebut. Keadaan leluhur dimiliki oleh nenek moyang seluruh kelompok,
sehingga setiap \omterm{def:g10:biodiversity-scales:species}{spesies} dapat saja mempertahankannya: berbagi keadaan
itu tidak menunjukkan apa-apa tentang kekerabatan yang lebih dekat.
\end{solution}

\begin{solution}{exo:g12:phylogenetic-trees:4}
Seorang nenek moyang dan semua keturunannya. Burung adalah sebuah
\omterm{prop:g12:phylogenetic-trees:rules}{klad}. ``Reptil'' dalam arti sehari-hari bukan \omterm{prop:g12:phylogenetic-trees:rules}{klad}: kelompok itu
meninggalkan burung, yang justru turun dari dalamnya.
\end{solution}

\begin{solution}{exo:g12:phylogenetic-trees:5}
Simpanse; simpul tertua adalah simpul yang menyatukan makaka dengan
kera, 25 juta tahun lalu.
\end{solution}

\begin{solution}{exo:g12:phylogenetic-trees:6}
Putarlah pada setiap simpul: akar terbelah menjadi katak (di bawah)
dan sisanya; lalu kadal, lalu mencit dan manusia di atas. Mencit dan
manusia tetap berbagi simpul termuda, kadal bergabung berikutnya, dan
katak paling akhir.
\end{solution}

\begin{solution}{exo:g12:phylogenetic-trees:7}
Merpati punya rahang, tungkai dan amnion: ia masuk \omterm{prop:g12:phylogenetic-trees:rules}{klad} amniota di
samping kadal dan mencit. Bulu, yang hanya ada pada satu \omterm{def:g10:biodiversity-scales:species}{spesies}, tidak
menentukan kelompok apa pun di sini; bulu akan menentukan burung bila
lebih banyak burung ditambahkan.
\end{solution}

\begin{solution}{exo:g12:phylogenetic-trees:8}
Rambut, susu dan rahangnya menempatkan lumba-lumba di antara mamalia;
menempatkannya bersama hiu akan menuntut semua karakter mamalia muncul
dua kali. Sirip dan tubuh ramping itu muncul secara terpisah pada
kedua garis keturunan: sebuah konvergensi.
\end{solution}

\begin{solution}{exo:g12:phylogenetic-trees:9}
$2.2/0.27 \approx 8$ juta tahun.
\end{solution}

\begin{solution}{exo:g12:phylogenetic-trees:10}
Keduanya adalah ujung masa kini; tak satu pun turun dari yang lain.
Yang benar: manusia dan simpanse adalah \omterm{def:g10:biodiversity-scales:species}{spesies} bersaudari, keturunan
satu nenek moyang bersama, kini punah, yang hidup sekitar enam juta
tahun lalu dan bukan manusia maupun simpanse.
\end{solution}

\begin{solution}{exo:g12:phylogenetic-trees:11}
Lamprei diketahui terletak di luar kelompok \omterm{def:g10:common-ancestry:bodyplan}{vertebrata} berahang,
sehingga keadaannya bersifat leluhur bagi mereka. Dengan mencit
sebagai kelompok luar, rambut, amnion, tungkai dan rahang semuanya
akan terbaca sebagai leluhur dan ketiadaannya sebagai turunan: pohonnya
akan terbangun terbalik.
\end{solution}

\begin{solution}{exo:g12:phylogenetic-trees:12}
\omterm{def:g12:selection-drift-speciation:population}{Populasi} leluhurnya membawa beberapa \omterm{def:g10:universal-dna:gene}{alel} tiap \omterm{def:g10:universal-dna:gene}{gen}; ketika ia terbelah
dua kali dalam waktu berdekatan, \omterm{def:g10:universal-dna:gene}{alel} yang berbeda terpilah ke dalam
ketiga keturunannya pada \omterm{def:g10:universal-dna:gene}{gen} yang berbeda, sehingga sejarah satu \omterm{def:g10:universal-dna:gene}{gen}
mengelompokkan \omterm{def:g10:biodiversity-scales:species}{spesies} 1 dengan 2 dan sejarah \omterm{def:g10:universal-dna:gene}{gen} lain mengelompokkan
2 dengan 3. Setiap pohon \omterm{def:g10:universal-dna:gene}{gen} benar bagi \omterm{def:g10:universal-dna:gene}{gennya} sendiri; pohon \omterm{def:g10:biodiversity-scales:species}{spesies}
adalah putusan mayoritas banyak \omterm{def:g10:universal-dna:gene}{gen}.
\end{solution}

\begin{solution}{exo:g12:phylogenetic-trees:13}
Pada 1\% per tahun, cuplikan virus yang diambil berselang beberapa
bulan sudah berbeda secara terukur: pohon virus itu menguraikan siapa
menulari siapa di dalam sebuah wabah, tetapi sepanjang jutaan tahun
urutannya akan tertimpa berkali-kali. Pada 0.3\% per juta tahun, \omterm{def:g10:universal-dna:gene}{gen}
mamalia berubah terlalu lambat untuk memperlihatkan apa pun di dalam
sebuah wabah tetapi menyimpan rekaman yang terbaca sepanjang 200 juta
tahun.
\end{solution}

\begin{solution}{exo:g12:phylogenetic-trees:14}
Fosil adalah \omterm{def:g10:biodiversity-scales:species}{spesies} tersendiri, dengan keadaan turunannya sendiri,
dan belum tentu \omterm{def:g12:selection-drift-speciation:population}{populasi} leluhurnya; ia duduk pada cabang samping
dekat simpul tempat bulu muncul. Fosil itu menunjukkan bahwa simpulnya
--- nenek moyang bersama burung dan kerabat reptilnya yang terdekat
--- berumur sekurang-kurangnya 150 juta tahun dan bahwa bulu mendahului
hilangnya gigi dan ekor.
\end{solution}

\begin{solution}{exo:g12:phylogenetic-trees:15}
Sebuah pohon tak punya puncak: memutar simpul mana pun memindahkan
manusia ke bawah tanpa mengubah satu pun kekerabatan. Setiap ujung
masa kini telah berevolusi selama waktu yang sama sejak akarnya,
sehingga tak satu pun lebih maju. Dan manusia adalah satu ujung di
antara jutaan, pada satu cabang kera; bentuk pohon merekam penurunan,
bukan peringkat.
\end{solution}

\begin{solution}{pb:g12:phylogenetic-trees:1}
\textbf{1.} Rangka bertulang (lima \omterm{def:g10:biodiversity-scales:species}{spesies}): \omterm{def:g10:common-ancestry:bodyplan}{vertebrata} bertulang;
hiu berada di luar.

\textbf{2.} Empat tungkai: tetrapoda (salamander, buaya, merpati,
kucing). Telur amniota: amniota (buaya, merpati, kucing). Bersarang:
\omterm{def:g10:common-ancestry:bodyplan}{vertebrata} bertulang $\supset$ tetrapoda $\supset$ amniota.

\textbf{3.} Tidak: keadaan pada satu \omterm{def:g10:biodiversity-scales:species}{spesies} saja tidak
mengelompokkan apa pun. Ketiganya akan berguna bila lebih banyak
\omterm{def:g10:biodiversity-scales:species}{spesies} berbagi keadaan itu (burung lain, mamalia lain, reptil lain).

\textbf{4.} Akar: hiu lawan sisanya (rangka bertulang); lalu salmon
lawan tetrapoda (empat tungkai); lalu salamander lawan amniota (telur
amniota); di antara amniota, ada percabangan bertiga yang belum
terurai, dengan bulu pada cabang merpati serta rambut dan lubang
tengkorak pada cabang kucing.

\textbf{5.} Tabel itu menempatkan buaya, merpati dan kucing bersama
tetapi tak dapat mengurutkannya: kerabat terdekat buaya belum putus.

\textbf{6.} Buaya--merpati, 8\%.

\textbf{7.} Kucing (15\% dari masing-masing); lalu salamander
(19--20\%); lalu salmon (24--25\%); lalu hiu (30--31\%).

\textbf{8.} Keduanya sepakat pada setiap sarang; \omterm{def:g10:universal-dna:information}{DNA} menguraikan
amniota: buaya dan merpati bersaudari, kucing bergabung berikutnya.

\textbf{9.} Garis keturunan hiu berpisah sebelum semua yang lain
memencar satu sama lain; masing-masing telah terpisah dari hiu selama
waktu yang sama dan menunjukkan jarak yang sama.

\textbf{10.} Buaya lebih dekat kerabatnya kepada merpati daripada
kepada kucing --- dan, dengan penalaran yang sama, daripada kepada
kadal; ``reptil'' tanpa burung bukanlah sebuah \omterm{prop:g12:phylogenetic-trees:rules}{klad}.

\textbf{11.} 19\% sepanjang $2 \times 340$ juta tahun cabang: sekitar
0.028\% per juta tahun per garis keturunan, yaitu 0.056\% per juta
tahun pemisahan.

\textbf{12.} Buaya--merpati $8/0.056 \approx 140$ juta tahun; simpul
amniota $15/0.056 \approx 270$ juta tahun.

\textbf{13.} Salmon $24.5/0.056 \approx 440$; hiu $30.5/0.056
\approx 540$ juta tahun.

\textbf{14.} Fosil: 250 dan 310; jam itu memberi 140 dan 270.
Urutannya benar dan tanggal yang lebih tua cukup dekat, tetapi tanggal
buaya--burung jauh terlalu muda: laju satu \omterm{def:g10:universal-dna:gene}{gen} berbeda antargaris
keturunan, dan satu \omterm{def:g10:universal-dna:gene}{gen} paling banter memberi jam yang kasar.

\textbf{15.} Pada jarak yang besar banyak posisi telah berubah lebih
dari sekali, sehingga perbedaannya tak lagi menjumlah: jam itu jenuh.
Simpul tua ditaksir terlalu muda oleh \omterm{def:g10:universal-dna:gene}{gen} ini; diperlukan \omterm{def:g10:universal-dna:gene}{gen} yang
lebih lambat, atau banyak \omterm{def:g10:universal-dna:gene}{gen} sekaligus.

\textbf{16.} (kucing, (buaya, merpati)) adalah \omterm{prop:g12:phylogenetic-trees:rules}{klad} terkecil yang
memuat kucing: amniota; lalu tetrapoda; lalu \omterm{def:g10:common-ancestry:bodyplan}{vertebrata} bertulang;
lalu keenamnya bersama hiu.

\textbf{17.} ``Ikan'' bukan sebuah \omterm{prop:g12:phylogenetic-trees:rules}{klad}: simpul yang menyatukan hiu
dan salmon adalah akarnya, yang keturunannya mencakup segalanya.
Amniota adalah \omterm{prop:g12:phylogenetic-trees:rules}{klad}: buaya, merpati, kucing dan nenek moyang bersama
mereka.

\textbf{18.} Salamander adalah \omterm{def:g10:biodiversity-scales:species}{spesies} masa kini, bukan perantara; ia
adalah kelompok saudari amniota, berbagi nenek moyang tetrapoda dengan
mereka dan tak punya telur amniota yang muncul kemudian pada garis
keturunan amniota.

\textbf{19.} Di samping cabang antara simpul tetrapoda dan simpul
amniota: seekor tetrapoda yang bukan amniota. Tarikhnya menegaskan
bahwa empat tungkai sudah ada pada 360 juta tahun lalu dan telur
amniota muncul kemudian.

\textbf{20.} (hiu, (salmon, (salamander, (kucing, (buaya,
merpati))))); \omterm{def:g10:universal-dna:information}{DNA} memutuskan bahwa buaya dan merpati bersaudari;
pemisahan keduanya, menurut \omterm{def:g10:universal-dna:gene}{gen} ini, bertarikh sekitar 140 juta tahun
(fosil menyebut 250).
\end{solution}
